\documentclass[11pt,a4paper]{article}
\usepackage{curso}
\begin{document}
\portada{Ejercicios --- Lección}{9.2}{Difusión discreta}
\section{Ejercicios}
\begin{enumerate}[label=\textbf{E\arabic*.}]
\item (B) Probabilidad de que un token esté enmascarado en $t$. ¿Es markoviano el proceso?
\item (B) ¿Por qué revelar un token por paso es exacto con un modelo perfecto?
\item (B) Con dos bits perfectamente correlacionados ($00$ o $11$, cada uno con 1/2), ¿qué produce un muestreador de un paso?
\item (B) Nombre dos formas de reducir los pasos.
\item (I) Para la cadena de la lección, calcule la probabilidad de las dos cadenas constantes.
\item (I) Demuestre que el muestreador de un paso produce cada bit con marginal $1/2$ y bits independientes.
\item (A) Implemente revelado por confianza (revelar primero las posiciones con condicional más extrema) y mida la variación total con 2 pasos.
\end{enumerate}
\section{Soluciones}
\begin{solucion}
\textbf{E1.} $t$; sí, con la máscara absorbente la probabilidad crece de forma monótona.

\textbf{E2.} Es la regla de la cadena: $p(x)=\prod_kp(x_{\sigma_k}\mid x_{\sigma_1},\dots,x_{\sigma_{k-1}})$.

\textbf{E3.} Bits independientes con marginal $1/2$: $00,01,10,11$ con $1/4$; variación total $1/2$.

\textbf{E4.} Revelado por confianza; destilación (o muestreo conjunto con un EBM).

\textbf{E5.} $2\times\frac12\times0{,}9^7=0{,}478$.

\textbf{E6.} Sin tokens visibles, cada condicional es la marginal ($1/2$ por simetría) y se muestrean por separado.

\textbf{E7.} Compare con el orden aleatorio (variación total 0,46 con 2 pasos); el resultado depende de la estructura de correlaciones.
\end{solucion}
\end{document}
